(1) 取 $H=H_0+\lambda x^4$，无扰动本征能量 $E_n^{(0)}=\hbar\omega(n+1/2)$。非简并Rayleigh–Schrödinger微扰论的一级公式为 $\Delta E_n^{(1)}=\lambda\langle n|x^4|n\rangle$，要求扰动能量相对于附近未扰动能级间隔足够小。
(2) 令 $x_0=\sqrt{\hbar/(2m\omega)}$，由 $x=x_0(a+a^\dagger)$ 及 $a|n\rangle=\sqrt n|n-1\rangle$、$a^\dagger|n\rangle=\sqrt{n+1}|n+1\rangle$，可得 $\langle n|(a+a^\dagger)^4|n\rangle=6n^2+6n+3$。故 $\langle n|x^4|n\rangle=3(2n^2+2n+1)x_0^4$，含正确的长度四次量纲。
(3) 一级修正 $\Delta E_n^{(1)}=3\lambda(2n^2+2n+1)\hbar^2/(4m^2\omega^2)$，基态修正为 $3\lambda\hbar^2/(4m^2\omega^2)$。因 $n=1$ 与 $n=0$ 的括号分别为5与1，首个跃迁能量差 $E_1-E_0=\hbar\omega+3\lambda\hbar^2/(m^2\omega^2)+O(\lambda^2)$，故跃迁角频率为 $\omega_{10}=\omega+3\lambda\hbar/(m^2\omega^2)+O(\lambda^2)$。
(4) 在低激发态可用无量纲小量 $\epsilon=\lambda\hbar/(m^2\omega^3)\ll1$ 衡量微扰强度；较高数态还需随 $n$ 增大收紧条件。$\lambda\to0$ 时修正消失、能级恢复等间距 $\hbar\omega$。绝对能量可以加上共同的任意参考常数，而跃迁频率来自可测能级差，因此更能直接体现非谐修正。
